\label{sec:concl}
\textbf{What we found.} (1) The released ballbot pipeline did not train on a stock simulator for three independent reasons (wrong friction frames on two of three wheels, a velocity-indexing bug that made the directional reward pay for angular velocity, and terrain 4--10$\times$ steeper than the tilt limit allows). After repairing them and re-tuning PPO (100\,Hz control, $\gamma{=}0.999$, a speed-capped reward and a fall penalty), policies learn to balance within $\approx$0.25M decisions and to drive at 0.1--0.2\,m/s on terrains of up to $\approx$9$^\circ$ mean slope, on held-out seeds. (2) \textbf{H1 is not supported}: frozen-encoder depth features did not lower the fall rate (they raised it on the steepest terrain) and gave no reliable success benefit once per-seed variation is considered. (3) \textbf{H2 is not observed} because the policies are insensitive even to fully blank depth, i.e.\ they ignore vision. (4) \textbf{H3 is inconclusive}: neither waypoint policy reached the goal in 8\,s.

\textbf{Limitations.} Two training seeds per variant (the pooled tests treat terrains as units and are dominated by seed variance); a 1M-decision budget, about one eighth of the baseline's 8M steps; a single frozen encoder pretrained on a different terrain distribution; success defined as 1\,m of progress because the policies are slow; one waypoint run; simulation only. The Python contact-frame fix reproduces the published patch's intent, not its exact numerical behaviour, and the NumPy gradient-noise fallback changes the exact terrains. Our absolute returns are not comparable to \cite{salehi2025} (different control rate and reward weights).

\textbf{Next steps.} (a) Make depth useful: train the encoder end-to-end or fine-tune it on this terrain distribution (\texttt{depth\_e2e} is implemented), add terrain types with features that cannot be felt before contact (steps, ridges), and use a harder curriculum so that proprioception alone is insufficient. (b) More seeds ($\ge5$) and a longer budget, with checkpoint selection on held-out terrains. (c) Re-run H3 with longer episodes and distance-based reward shaping. (d) Richer sensors and planning: a learned or map-based waypoint planner on top of the low-level policy (hierarchical control), and sim-to-real considerations such as domain randomisation.
